% !PRINTED 677
\weberpage{677}
%<<CONT
darstellen, worin $\mathfrak{U}$ und $\mathfrak{U}'$ äquivalente $n$-Ecke ohne gemeinschaftlichen Teiler sind. Eine ganze Funktion $\omega$ von $z$ ist dadurch charakterisiert, daß sie in keinem Punkt unendlich wird, der nicht in $\mathfrak{U}$ enthalten ist. Folglich läßt sich eine solche Funktion $\omega$ so darstellen:
\weq{2}{\omega = \dfrac{\mathfrak{A}}{\mathfrak{U}^r},}
worin $r$ ein positiver Exponent ist. In (2) können Zähler und Nenner gemeinschaftliche Teiler haben. Wir können aber annehmen, daß $\mathfrak{A}$ nicht durch $\mathfrak{U}$ teilbar ist. Dann hat in (2) der Exponent $r$ den kleinst möglichen Wert. \gesperrt{Diese Zahl $r$ soll der Exponent von $\omega$ in $z$ heißen.}

1. \gesperrt{Setzen wir $z' = 1:z$, so ist nach (1), (2):}
\[
z' = \frac{\mathfrak{U}}{\mathfrak{U}'}, \quad \omega' = \omega z'^r = \frac{\mathfrak{A}}{\mathfrak{U}'^r},
\]
und wenn $\mathfrak{A}$ nicht durch $\mathfrak{U}'$ teilbar ist, so ist $r$ der Exponent von $\omega'$ in bezug auf $z'$. Die Annahme, daß $\mathfrak{A}$ nicht durch $\mathfrak{U}'$ teilbar sei, bedeutet, daß $\omega$ nicht durch $z$ teilbar sei, und daraus folgt, daß auch $\omega'$ nicht durch $z'$ teilbar ist.

2. Der Exponent einer ganzen rationalen Funktion von $z$ vom $m$ten Grade
\[
x = a_0 + a_1 z + \cdots + a_m z^m
\]
ist gleich $m$. Denn diese Funktion wird in jedem in $\mathfrak{U}$ enthaltenen Punkt unendlich von der $m$ten Ordnung. Der Exponent des Produktes $x\omega$ ist gleich $r + m$.

3. Der Exponent einer Summe
\[
\omega = c_1\omega_1 + c_2\omega_2 + \cdots + c_s\omega_s
\]
ist höchstens gleich dem größten der Exponenten von $\omega_1, \omega_2, \ldots, \omega_s$.

4. Wir bestimmen eine Reihe ganzer Funktionen von $z$
\[
\lambda_1, \lambda_2, \lambda_3, \ldots
\]
durch folgende recurrente Bestimmung:
\begin{itemize}
\item[1)] $\lambda_1$ konstant (Exponent $r_1 = 0$),
\item[2)] $\lambda_2$ eine nicht rationale ganze Funktion von $z$ mit möglichst kleinem Exponenten $r_2$,
\item[3)] $\lambda_s$ eine nicht in der Form
\weq{3}{x_1\lambda_1 + x_2\lambda_2 + \cdots + x_{s-1}\lambda_{s-1}}
enthaltene ganze Funktion von $z$ mit möglichst kleinem Exponenten $r_s$,
\end{itemize}
%CONT>>
